% Source: https://artofproblemsolving.com/community/c3h2170502_angle_chasing
% Stable identity: aops-topic-2170502-q01
Let $P$ be the midpoint of the diagonal $BD$.

\begin{asy}
usepackage("amsmath");
size(6.5cm);
defaultpen(fontsize(10pt));
pair B=(0,0),C=(3,0); real L=2.2;
pair A=B+L*dir(150), D=C+L*dir(90);
pair M=(B+C)/2, Nn=(A+D)/2, P=(B+D)/2;
draw(A--B--C--D--cycle,linewidth(0.9));
draw(B--D,gray+dashed); draw(M--P--Nn,blue); draw(M--Nn,red+linewidth(1));
draw(C+(-0.2,0)--C+(-0.2,0.2)--C+(0,0.2));
dot(M);dot(Nn);dot(P);
label("$A$",A,W);label("$B$",B,SW);label("$C$",C,SE);label("$D$",D,NE);
label("$M$",M,S);label("$N$",Nn,N);label("$P$",P,E);
label("$150^\circ$",B+0.5*dir(108),fontsize(8pt));
\end{asy}

In triangle $BCD$, $MP$ joins the midpoints of $BC$ and $BD$, so $MP\parallel CD$ and $MP=\dfrac{1}{2}CD$. In triangle $DAB$, $PN$ joins the midpoints of $DB$ and $DA$, so $PN\parallel BA$ and $PN=\dfrac{1}{2}AB$. Since $AB=CD$, triangle $MPN$ is isosceles with $PM=PN$.

Direct $BC$ to the right. Then the ray $PM$ points straight down (opposite to $\overrightarrow{CD}$, which is perpendicular to $BC$), and the ray $PN$ has the direction of $\overrightarrow{BA}$, which makes $150^\circ$ with $\overrightarrow{BC}$. The direction straight down makes $270^\circ$, so $\angle MPN=270^\circ-150^\circ=120^\circ$, and the base angles of the isosceles triangle are $\angle PMN=\angle PNM=30^\circ$.

As $MP\perp BC$, the line $MN$ makes an angle of $90^\circ-30^\circ=\boxed{60^\circ}$ with $BC$.

(Without $AB=CD$ the angle is not determined: for $AB=2$, $BC=1$, $CD=\dfrac{1}{2}$ the acute angle is about $40.9^\circ$.)
